\documentclass{article}
\usepackage{amsmath}
\usepackage{propositions}

%% A bare string in an item's optional argument means name={...}.  This must
%% work whether or not the argument also carries genuine key=value entries.

\begin{document}

\section{Bare name alone}

\begin{prop}
  \item[No Overlap] Name with a space. \label{b:space}
  \item[Nihilism] Single word. \label{b:word}
  \item[$P \wedge Q$] Math, catcodes preserved. \label{b:math}
  \item[{Alpha, Beta}] Braced, containing a comma. \label{b:comma}
\end{prop}

References: \ref{b:space}, \ref{b:word}, \ref{b:math}, \ref{b:comma}.

\section{Bare name alongside other keys}

Each of these must show its name, not fall back to a number.

\begin{prop}
  \item[No Overlap, align=flush, display format=\textsc{#1},
    ref format=\textit{#1}]
    Bare name first, then keys with \texttt{\#1} in their values.
    \label{b:keys}
  \item[align=flush, Trailing Name] Bare name last. \label{b:last}
  \item[Nihilism, style=thesis] Bare name with \texttt{style}. \label{b:style}
  \item[{Alpha, Beta}, align=flush] Braced comma name with a key.
    \label{b:comma2}
  \item[Ext\emph{dash}, gloss=boo] A macro (with braces) in the name.
    \label{b:macro}
  \item[\ref{b:word}$'$, align=flush] A name built from another reference.
    \label{b:built}
\end{prop}

References: \ref{b:keys}, \ref{b:last}, \ref{b:style}, \ref{b:comma2},
\ref{b:macro}, \ref{b:built}.

\section{Keys that are not names}

\begin{prop}
  \item[reset] A valueless key must stay a key, so this item is numbered.
    \label{b:reset}
  \item[reset=false] Explicit form, also numbered. \label{b:reset2}
  \item[counter=none] Text item: no number, no name.  A counter does not make
    an item named, so this one takes the nameless style --- but with nothing to
    display it must produce no label at all, not that style's decoration
    wrapped around nothing (\texttt{()}).
  \item[counter=none, gloss=g] Text item with a gloss: the gloss survives.
  \item Plain numbered item. \label{b:plain}
\end{prop}

References: \ref{b:reset}, \ref{b:reset2}, \ref{b:plain}.

\section{Other entry points}

\begin{equation}
  x = y \ptag[Tag Me, align=flush]
\end{equation}

\begin{prop}
  \item[style=thesis, Shortcut Name, align=flush] Style shortcut command. \label{b:short}
\end{prop}

Reference: \ref{b:short}.

\end{document}
